\documentclass[a4paper,11pt,titlepage]{article}
\usepackage[utf8]{inputenc}
\usepackage[T1]{fontenc}
\usepackage{xeCJK}
\setCJKmainfont{Hiragino Kaku Gothic Pro}
\usepackage{geometry}
\geometry{left=25mm,right=25mm,top=30mm,bottom=30mm}
\usepackage{graphicx}
\usepackage{xcolor}
\usepackage{tikz}
\usepackage{tcolorbox}
\usepackage{booktabs}
\usepackage{longtable}
\usepackage{enumitem}
\usepackage{hyperref}
\usepackage{fancyhdr}
\usepackage{titlesec}

% Colors
\definecolor{primary}{RGB}{0,102,204}
\definecolor{secondary}{RGB}{51,51,51}
\definecolor{accent}{RGB}{255,153,0}
\definecolor{success}{RGB}{40,167,69}
\definecolor{danger}{RGB}{220,53,69}
\definecolor{lightgray}{RGB}{248,249,250}

% Header/Footer
\pagestyle{fancy}
\fancyhf{}
\fancyhead[L]{\textcolor{primary}{\small 安全書類管理システム提案書}}
\fancyhead[R]{\textcolor{secondary}{\small \thepage}}
\fancyfoot[C]{\textcolor{secondary}{\small Powered by Lark Base}}
\renewcommand{\headrulewidth}{0.5pt}

% Section styling
\titleformat{\section}{\Large\bfseries\color{primary}}{\thesection}{1em}{}
\titleformat{\subsection}{\large\bfseries\color{secondary}}{\thesubsection}{1em}{}

% Custom boxes
\newtcolorbox{infobox}{colback=lightgray,colframe=primary,title=}
\newtcolorbox{alertbox}{colback=danger!10,colframe=danger,title=}
\newtcolorbox{successbox}{colback=success!10,colframe=success,title=}

\begin{document}

% Title Page
\begin{titlepage}
\centering
\vspace*{2cm}

{\Huge\bfseries\textcolor{primary}{建設業向け\\安全書類管理システム}}

\vspace{1cm}

{\Large\textcolor{secondary}{Lark Base 完全自動化提案書}}

\vspace{2cm}

\begin{tikzpicture}
\draw[primary,line width=2pt] (0,0) -- (10,0);
\end{tikzpicture}

\vspace{2cm}

{\large
\begin{tabular}{rl}
\textbf{システム構成} & 37テーブル / 29自動化ビュー \\
\textbf{対象業務} & 10カテゴリ \\
\textbf{自動化率} & 80\%以上 \\
\end{tabular}
}

\vspace{3cm}

{\Large\textcolor{secondary}{作成日: 2024年11月21日}}

\vfill

{\large\textcolor{primary}{Miyabi AI Solutions}}

\end{titlepage}

% Table of Contents
\tableofcontents
\newpage

% Executive Summary
\section{エグゼクティブサマリー}

\begin{successbox}
\textbf{本提案の価値}\\
建設現場の安全書類管理業務を80\%以上自動化し、年間約500時間の工数削減を実現します。
\end{successbox}

\subsection{導入効果}

\begin{center}
\begin{tabular}{lcc}
\toprule
\textbf{項目} & \textbf{現状} & \textbf{導入後} \\
\midrule
書類作成時間 & 2時間/日 & 30分/日 \\
承認プロセス & 3日 & 即時〜1日 \\
検索・照会 & 30分 & 即時 \\
ミス発生率 & 5\% & 0.5\% \\
\bottomrule
\end{tabular}
\end{center}

\subsection{システム概要}

本システムは以下の10カテゴリで構成されています：

\begin{enumerate}[leftmargin=*]
\item \textbf{マスタ管理} - 会社・作業員・資格等の基本情報
\item \textbf{工事管理} - 現場・工期・進捗管理
\item \textbf{書類管理} - 提出物・持込機械届
\item \textbf{リース発注} - 機材発注から請求書まで自動化
\item \textbf{経費精算} - AI勘定科目判定・自動承認
\item \textbf{勤怠管理} - アルコールチェック・作業日報
\item \textbf{点検管理} - 始業前点検チェックリスト
\item \textbf{労務調査} - 労働時間・賃金管理
\item \textbf{車両管理} - 車検・燃費・運転日報
\item \textbf{安全教育} - 教育記録・事故報告・KY活動
\end{enumerate}

\newpage

% System Architecture
\section{システム構成}

\subsection{テーブル一覧（37テーブル）}

\subsubsection{マスタ系（12テーブル）}

\begin{longtable}{clp{8cm}}
\toprule
\textbf{No} & \textbf{テーブル名} & \textbf{用途} \\
\midrule
\endhead
M1 & 元請け会社 & 発注元企業の情報管理 \\
M2 & 下請け・協力会社 & 施工会社・建設業許可管理 \\
M3 & 資格マスタ & 国家資格・技能講習の定義 \\
M4 & 安全書類テンプレート & 提出書類の種類と要件 \\
M5 & 職人・作業員 & 作業員の個人情報・資格 \\
M6 & 勘定科目マスタ & 経費分類・自動振り分けキーワード \\
M7 & 点検テンプレート & 点検種別・法定点検の定義 \\
M8 & 点検項目マスタ & チェック項目・判定基準 \\
M9 & 安全教育マスタ & 教育種別・法的根拠 \\
M10 & リース会社 & リース業者情報 \\
M11 & リース品目マスタ & レンタル機材の単価・分類 \\
M12 & 車両マスタ & 社用車・車検・保険管理 \\
\bottomrule
\end{longtable}

\subsubsection{トランザクション系（25テーブル）}

各業務の日々の記録を管理するテーブル群です。

\begin{center}
\begin{tabular}{lc}
\toprule
\textbf{カテゴリ} & \textbf{テーブル数} \\
\midrule
工事管理 & 1 \\
書類管理 & 2 \\
リース発注 & 5 \\
経費精算 & 3 \\
勤怠管理 & 2 \\
点検管理 & 2 \\
労務調査 & 3 \\
車両管理 & 3 \\
安全教育 & 4 \\
\bottomrule
\end{tabular}
\end{center}

\newpage

% Business Flows
\section{業務フロー}

\subsection{フロー1: リース発注〜請求書申請}

\begin{tcolorbox}[colback=lightgray,colframe=primary]
\textbf{自動化ポイント}: テンプレートから一括発注生成、請求書の自動マッチング
\end{tcolorbox}

\begin{enumerate}
\item \textbf{年度計画作成}
    \begin{itemize}
    \item 年間発注テンプレートに予定品目・数量を登録
    \item 昨年実績を参照して調整
    \end{itemize}

\item \textbf{発注書自動生成}
    \begin{itemize}
    \item テンプレートから発注書を一括生成
    \item 発注番号は自動採番（PO-YYYYMMDD-001）
    \end{itemize}

\item \textbf{リース会社へ送信}
    \begin{itemize}
    \item 発注方法（メール/FAX）を選択
    \item 送信日時を自動記録
    \end{itemize}

\item \textbf{納品登録}
    \begin{itemize}
    \item 発注明細の「納品済み」をチェック
    \item 持込機械テーブルに自動連携
    \end{itemize}

\item \textbf{請求書受領・申請}
    \begin{itemize}
    \item 請求書画像をアップロード
    \item 関連発注書を自動マッチング
    \item 承認ワークフロー開始
    \end{itemize}
\end{enumerate}

\subsection{フロー2: 経費精算（AI自動分類）}

\begin{tcolorbox}[colback=lightgray,colframe=primary]
\textbf{自動化ポイント}: レシート画像から勘定科目をAI判定、10,000円以上は自動で承認リクエスト
\end{tcolorbox}

\begin{enumerate}
\item \textbf{レシート撮影・アップロード}
    \begin{itemize}
    \item 現場監督がスマホで撮影
    \item Larkアプリからダイレクト投稿
    \end{itemize}

\item \textbf{AI勘定科目判定}
    \begin{itemize}
    \item 店舗名・摘要からキーワードマッチング
    \item 勘定科目を自動設定
    \item 信頼度スコアを表示
    \end{itemize}

\item \textbf{承認プロセス}
    \begin{itemize}
    \item 10,000円未満: 自動承認
    \item 10,000円以上: 承認者に通知
    \end{itemize}

\item \textbf{月次レポート自動集計}
    \begin{itemize}
    \item 現場別・科目別の集計
    \item 予算対比グラフ
    \end{itemize}
\end{enumerate}

\newpage

\subsection{フロー3: 勤怠管理・アルコールチェック}

\begin{alertbox}
\textbf{法的要件}: 2022年4月より、白ナンバー事業者もアルコールチェック記録が義務化
\end{alertbox}

\begin{enumerate}
\item \textbf{出勤打刻}
    \begin{itemize}
    \item 現場到着時にLarkアプリで打刻
    \item GPS位置情報を自動記録
    \end{itemize}

\item \textbf{アルコールチェック実施}
    \begin{itemize}
    \item 検知器の数値を入力
    \item 異常検出時は即座に管理者通知
    \item 「🚫検出」フラグが自動設定
    \end{itemize}

\item \textbf{体調確認}
    \begin{itemize}
    \item 体温・体調を記録
    \item 不調時は別作業への配置換え
    \end{itemize}

\item \textbf{退勤打刻・日報}
    \begin{itemize}
    \item 退勤時刻を記録
    \item 作業日報に当日の作業内容を入力
    \end{itemize}
\end{enumerate}

\subsection{フロー4: 始業前点検}

\begin{tcolorbox}[colback=lightgray,colframe=primary]
\textbf{自動化ポイント}: チェックリストをスマホで完結、不具合は即時通知
\end{tcolorbox}

\begin{enumerate}
\item \textbf{点検テンプレート選択}
    \begin{itemize}
    \item 機械種別に応じたテンプレートを選択
    \item バックホウ、足場、高所作業車等
    \end{itemize}

\item \textbf{チェック項目入力}
    \begin{itemize}
    \item 各項目を「良好/要注意/不良」で評価
    \item 必須項目は写真添付
    \end{itemize}

\item \textbf{結果判定}
    \begin{itemize}
    \item 全項目良好 → 「合格」
    \item 不良あり → 「要修理」「使用禁止」
    \end{itemize}

\item \textbf{是正措置}
    \begin{itemize}
    \item 不具合内容を記録
    \item 是正措置を実施
    \item 完了日を記録
    \end{itemize}
\end{enumerate}

\newpage

\subsection{フロー5: 安全教育・KY活動}

\begin{enumerate}
\item \textbf{教育計画}
    \begin{itemize}
    \item 月次の教育スケジュールを設定
    \item 新規入場者教育、職長教育等
    \end{itemize}

\item \textbf{教育実施・記録}
    \begin{itemize}
    \item 出席者を登録
    \item 教育内容・テスト結果を記録
    \item 修了証番号を発行
    \end{itemize}

\item \textbf{資格有効期限管理}
    \begin{itemize}
    \item 期限1ヶ月前に自動通知
    \item 再教育の手配
    \end{itemize}

\item \textbf{KY活動記録}
    \begin{itemize}
    \item 毎朝の危険予知活動を記録
    \item 危険項目・対策・行動目標
    \end{itemize}
\end{enumerate}

\subsection{フロー6: 事故・ヒヤリハット報告}

\begin{alertbox}
\textbf{重要}: ヒヤリハット300件で重大事故1件（ハインリッヒの法則）
\end{alertbox}

\begin{enumerate}
\item \textbf{即時報告}
    \begin{itemize}
    \item 発生直後にスマホから報告
    \item 現場写真を添付
    \end{itemize}

\item \textbf{重大度判定}
    \begin{itemize}
    \item Ⅰ（重大）→ 🔴緊急対応
    \item Ⅱ（中程度）→ 🟡要注意
    \item Ⅲ〜Ⅳ → 🟢通常対応
    \end{itemize}

\item \textbf{原因分析・再発防止}
    \begin{itemize}
    \item 原因を分析
    \item 再発防止策を策定
    \item 是正完了日を記録
    \end{itemize}

\item \textbf{月次分析}
    \begin{itemize}
    \item 発生傾向を分析
    \item 安全大会での共有
    \end{itemize}
\end{enumerate}

\newpage

\subsection{フロー7: 労務調査対応}

\begin{enumerate}
\item \textbf{月次労務データ収集}
    \begin{itemize}
    \item 作業員ごとの労働時間を集計
    \item 勤怠記録から自動連携
    \end{itemize}

\item \textbf{賃金計算}
    \begin{itemize}
    \item 賃金台帳の単価を適用
    \item 残業・深夜・休日手当を自動計算
    \end{itemize}

\item \textbf{労務報告書作成}
    \begin{itemize}
    \item 下請け会社が報告書を作成
    \item 元請け会社に提出
    \end{itemize}

\item \textbf{元請け集計・提出}
    \begin{itemize}
    \item 全下請けの報告書を集計
    \item 調査機関へ提出
    \end{itemize}
\end{enumerate}

\subsection{フロー8: 車両管理}

\begin{enumerate}
\item \textbf{日常管理}
    \begin{itemize}
    \item 運転日報で使用状況を記録
    \item 走行距離を自動累計
    \end{itemize}

\item \textbf{燃料管理}
    \begin{itemize}
    \item 給油記録を登録
    \item 燃費を自動計算
    \end{itemize}

\item \textbf{点検・車検管理}
    \begin{itemize}
    \item 車検残日数を自動表示
    \item 30日前に自動通知
    \end{itemize}

\item \textbf{保険管理}
    \begin{itemize}
    \item 自賠責・任意保険の期限管理
    \item 更新漏れ防止
    \end{itemize}
\end{enumerate}

\newpage

% Automation Features
\section{自動化機能}

\subsection{作成済みビュー（29ビュー）}

自動フィルタリングされたビューにより、即座に必要な情報にアクセスできます。

\begin{longtable}{lll}
\toprule
\textbf{テーブル} & \textbf{ビュー名} & \textbf{用途} \\
\midrule
\endhead
経費申請 & 📋未承認一覧 & 承認待ち案件の確認 \\
経費申請 & 🔴却下案件 & 再申請が必要な案件 \\
発注書 & 📦納品待ち & 未納品の追跡 \\
発注書 & 🗓️カレンダー & 納期のガントチャート \\
点検記録 & ⚠️要対応 & 不具合機械の確認 \\
点検記録 & 📊点検カンバン & 進捗のカンバン表示 \\
勤怠記録 & 🚫アルコール検出 & 異常検出者の確認 \\
請求書申請 & 🔴期限超過 & 支払い遅延案件 \\
請求書申請 & 🟡期限間近 & 7日以内の支払い \\
車両マスタ & ⚠️車検期限 & 30日以内の車検 \\
事故報告 & 🔴重大案件 & 緊急対応が必要 \\
提出物 & 📋未提出 & 未提出書類の確認 \\
提出物 & 📊進捗カンバン & 提出状況の可視化 \\
\bottomrule
\end{longtable}

\subsection{推奨Automation設定}

Lark Base のAutomation機能で以下を設定してください：

\begin{enumerate}
\item \textbf{期限アラート}
    \begin{itemize}
    \item トリガー: 毎日9:00
    \item 条件: 車検残日数 < 7
    \item アクション: 車両担当者にLark通知
    \end{itemize}

\item \textbf{承認リクエスト}
    \begin{itemize}
    \item トリガー: 経費申請レコード作成時
    \item 条件: 金額 >= 10,000
    \item アクション: 承認者にLark通知
    \end{itemize}

\item \textbf{アルコール異常検知}
    \begin{itemize}
    \item トリガー: 勤怠記録更新時
    \item 条件: アルコール検出値 > 0
    \item アクション: 管理者に即時通知
    \end{itemize}

\item \textbf{重大事故報告}
    \begin{itemize}
    \item トリガー: 事故報告レコード作成時
    \item 条件: 重大度 = Ⅰ
    \item アクション: 経営層に即時通知
    \end{itemize}
\end{enumerate}

\newpage

% Implementation Guide
\section{導入・運用ガイド}

\subsection{Phase 1: 初期設定（1週間）}

\begin{enumerate}
\item \textbf{マスタデータ登録}
    \begin{itemize}
    \item 元請け・下請け会社情報
    \item 作業員・資格情報
    \item 車両・機材情報
    \end{itemize}

\item \textbf{ユーザー権限設定}
    \begin{itemize}
    \item 管理者: 全テーブル編集可
    \item 現場監督: 担当現場のみ編集可
    \item 作業員: 入力フォームのみ
    \end{itemize}

\item \textbf{Automation設定}
    \begin{itemize}
    \item 通知ルールの設定
    \item 自動化ワークフローの有効化
    \end{itemize}
\end{enumerate}

\subsection{Phase 2: 試験運用（2週間）}

\begin{enumerate}
\item \textbf{パイロット現場で運用開始}
\item \textbf{フィードバック収集}
\item \textbf{設定調整}
\end{enumerate}

\subsection{Phase 3: 本格運用}

\begin{enumerate}
\item \textbf{全現場への展開}
\item \textbf{定期的な運用レビュー}
\item \textbf{継続的な改善}
\end{enumerate}

\newpage

% Benefits
\section{導入効果}

\subsection{定量的効果}

\begin{center}
\begin{tabular}{lrrr}
\toprule
\textbf{業務} & \textbf{現状時間} & \textbf{削減時間} & \textbf{削減率} \\
\midrule
経費精算 & 30分/件 & 5分/件 & 83\% \\
点検記録 & 20分/回 & 5分/回 & 75\% \\
日報作成 & 15分/日 & 3分/日 & 80\% \\
書類検索 & 30分/回 & 即時 & 100\% \\
月次集計 & 4時間/月 & 自動 & 100\% \\
\bottomrule
\end{tabular}
\end{center}

\begin{successbox}
\textbf{年間削減効果}\\
\begin{itemize}
\item 工数削減: 約500時間/年
\item コスト削減: 約150万円/年（人件費換算）
\item ペーパーレス: 約90\%削減
\end{itemize}
\end{successbox}

\subsection{定性的効果}

\begin{enumerate}
\item \textbf{コンプライアンス強化}
    \begin{itemize}
    \item 法定点検の確実な実施
    \item アルコールチェック記録の完備
    \item 監査対応の迅速化
    \end{itemize}

\item \textbf{安全性向上}
    \begin{itemize}
    \item ヒヤリハット報告の活性化
    \item 危険の早期発見
    \item 事故発生率の低減
    \end{itemize}

\item \textbf{業務効率化}
    \begin{itemize}
    \item リアルタイムな情報共有
    \item 承認プロセスの迅速化
    \item 二重入力の排除
    \end{itemize}
\end{enumerate}

\newpage

% Next Steps
\section{今後の拡張}

\subsection{追加機能案}

\begin{enumerate}
\item \textbf{モバイルアプリ最適化}
    \begin{itemize}
    \item 現場向けカスタムアプリ
    \item オフライン対応
    \end{itemize}

\item \textbf{AI機能強化}
    \begin{itemize}
    \item レシートOCR自動読取
    \item 点検写真の異常検知
    \end{itemize}

\item \textbf{外部連携}
    \begin{itemize}
    \item 会計ソフト連携
    \item 電子契約連携
    \end{itemize}

\item \textbf{ダッシュボード高度化}
    \begin{itemize}
    \item 経営指標KPI
    \item 安全スコアボード
    \end{itemize}
\end{enumerate}

\subsection{サポート体制}

\begin{itemize}
\item 導入支援: 初期設定・データ移行サポート
\item 運用支援: 問い合わせ対応・設定変更
\item 教育支援: 利用者向けトレーニング
\end{itemize}

\vspace{2cm}

\begin{center}
\begin{tcolorbox}[colback=primary!10,colframe=primary,width=12cm]
\centering
\Large\textbf{お問い合わせ}\\
\vspace{0.5cm}
\normalsize
Lark Base: \url{https://customer-cloud.jp.larksuite.com/base/HJDMbTdNjaJC8FsuVKejDOGJpLb}\\
\vspace{0.3cm}
Miyabi AI Solutions
\end{tcolorbox}
\end{center}

\end{document}
